\section{Holomorphic Function}

\subsection{Definitions}

\begin{definition}
    [Disconnection] If $E$ is disconnected, then there are $A,B$ nonempty such that
    \[
    E = A\cup B,\quad \overline{A} \cap B = A\cap \overline{B} = \varnothing
    \]
    If $E$ is closed/open and disconnected, then $E$ is a union of two disjoint closed/open set.
\end{definition}
\begin{proof}
    If $E$ is closed and $E = A \cup B$, then since $\overline{A} \cap B = \varnothing$, then $\overline{A} \subset E - B \subset A$ and hence $A$ closed. Similarly, $B$ closed.\par
    If $E$ is open with $A\cup B = E$, then $ B \subset E - \overline{A} \subset B$ and hence $B = E - \overline{A}$ open. Similarly, $A$ open and we are done.
\end{proof}

\begin{definition}
    [Holomorphic Function] For an open subset $\Omega$ of complex plane, and $f: \Omega \to \C$, for $z_0 \in \Omega$ and the limit
    \[
    \lim_{z\to z_0} \dfrac{f(z)-f(z_0)}{z-z_0}
    \]
    exists, denoted as $f'(z_0)$. If for any $z\in\Omega$, $f'(z)$ exists, then we say $f$ is holomorphic in $\Omega$. Denote $H(\Omega)$ as all the functions holomorphic in $\Omega$.
\end{definition}

\begin{definition}
    [Series] If for any $D(a;r) \subset \Omega$, there is a series
    \[
    \sum\limits_{n=0}^{\infty} c_n(z-a)^n
    \]
    which converges to $f(z)$ for all $z \in D(a;r)$, we say that $f$ has a series represenation in $\Omega$.
\end{definition}

\begin{theorem}
\label{thm: convergent-radius}
    For a series
    \[
    \sum\limits_{n=0}^{\infty} c_n(z-a)^n,
    \]
    for any $r < R$, we have the series is absolutely convergent in $\overline{D}(a;r)$ and for $z\notin \overline{D}(a;R)$, the series diverge where
    \[
    R = \dfrac{1}{\limsup_{n\to\infty} |c_n|^{1/n}}.
    \]
\end{theorem}
\begin{proof}
    For $r < R$, assume $r = q R, q < 1$ and then
    \[
    \sum\limits_{n=0}^{\infty} |c_n||z-a|^n \leq \sum\limits_{n = 0}^{\infty} q^n|c_n|(\limsup_{n\to\infty} |c_n|)^{-1} < \infty.
    \]
    For $|z-a| > R$, we have
    \[
    |c_n(z-a)^n| > |c_n|(\limsup_{n\to\infty} |c_n|)^{-1} \geq |c_n|(\sup_{m\geq N} |c_n|)^{-1}
    \]
    for $N < n$, which means we may always choose $n > N$ such that term above is larger than $1/2$ and we know the series can not converge absolutely since the series is not Cauchy.
\end{proof}

\begin{theorem}
    If $f$ has a series representation in $\Omega$, then $f\in H(\Omega)$ and $f'$ also has a series representation in $\Omega$. In particular, if for $z\in D(a;r) \subset \Omega$, we have
    \[
    f(z) = \sum\limits_{n=0}^{\infty} c_n(z-a)^n,
    \]
    then
    \[
    f'(z) = \sum\limits_{n=1}^{\infty}nc_n(z-a)^{n-1}.
    \]
\end{theorem}
\begin{proof}
    If in $D(a;r)$ we have
    \[\sum\limits_{n=0}^{\infty} c_n(z-a)^n\]
    to be convergent, then
    \[
    g(z): = \sum\limits_{n=1}^{\infty}nc_n(z-a)^{n-1}
    \]
    is also convergent in $D(a;r)$ by  the Theorem \ref{thm: convergent-radius}. And we have
    \[
    \begin{aligned}
        \lim_{z\to z_0} \left|\dfrac{f(z) - f(z_0)}{z-z_0} - g(z_0)\right| & = \lim_{z\to z_0}\left|\dfrac{1}{z-z_0}\left(\sum\limits_{n=0}^{\infty} c_n\left((z-a)^n-(z_0-a)^n\right)\right) -g(z)\right|\\
        &\leq \lim_{z\to z_0}\sum\limits_{n=0}^{\infty}|c_n|\left|\dfrac{(z-a)^n - (z_0-a)^n}{z-z_0} -n(z_0-a)^{n-1}\right| \\
        &\leq \lim_{z\to z_0}\sum\limits_{n=0}^{\infty}|c_n||z-z_0|\dfrac{n(n-1)}{2}r^{n-1} = 0
    \end{aligned}
    \]
    and we are done.
\end{proof}

\begin{corollary}
    If $f$ has a series representation in $\Omega$, then $f\in H(\Omega)$ and $f^{(k)}$ has a series representation for any integer $k$, in particular, $f^{(k)} \in H(\Omega)$ as well.
\end{corollary}

\begin{theorem}
\label{thm: route-integral-holomorphic}
    Assume $\mu$ is a complex finite measure on $X$, $\phi:X \to \C$ measurable and $\Omega$ is a disjoint open subset of $\C$ with $\phi(X)$, with
    \[
    f(z) = \int_X \dfrac{d\mu(\zeta)}{\phi(\zeta)- z}, \quad z\in \Omega
    \]
    then $f$ has a series represenation in $\Omega$.
\end{theorem}
\begin{proof}
    Assume that if $D(a;r) \subset \Omega$, then we know
    \[
    \left|\dfrac{z-a}{\phi(\zeta) - a}\right| \leq \left|\dfrac{z-a}{r}\right| < 1
    \]
    for any fixed $z\in D(a;r)$, and hence
    \[
    \dfrac{1}{1-\dfrac{z-a}{\phi(\zeta) - a}} = \dfrac{\phi(\zeta) -a }{\phi(\zeta) - z} 
    \]
    exists, which means 
    \[
    \sum\limits_{n=0}^{\infty}\dfrac{(z-a)^n}{(\phi(\zeta)-a)^{n+1}} = \dfrac{1}{\phi(\zeta) - z}.
    \]
    We may use DCT to see that
    \[
    f(z) = \sum\limits_{n=0}^{\infty} c_n(z-a)^n,\quad c_n:= \int_X \dfrac{d\mu(\zeta)}{(\phi(\zeta) -a)^{n+1}}.
    \]
\end{proof}

\subsection{Integral on Paths}

\begin{definition}
    [Path and Integral on Path] A curve in $X$ is a continuous map $\gamma:[a,b] \to X$ for some compact interval $[a,b]$.\par
    A path id s piecewisely differentiable curve $\gamma:[a,b] \to \C$ and the integral for $f$ continuous on $\gamma^*:= \{\gamma(t), t\in [a,b]\}$ is defined by
    \[
    \int_{\gamma}f(z)dz = \int_a^b f(\gamma(t))\gamma'(t)dt.
    \]
\end{definition}

\begin{lemma}
    If $\phi:[a,b] \to [a',b']$ is a continuously differentiable bijection such that $\phi(a) = a', \phi(b) = b'$, then for a path $\gamma:[a',b'] \to \C$ and $f$ continuous on $\gamma^*$, we have
    \[
    \int_{\gamma} f(z)dz = \int_{\gamma'}f(z)dz
    \]
    where
    \[\tilde{\gamma} := \gamma\circ \phi:[a,b] \to \C.\]
\end{lemma}
\begin{proof}
    By simple computation:
    \[
    \begin{aligned}
        \int_{\tilde{\gamma}}f(z)dz &= \int_a^b f(\tilde{\gamma}(t))(\tilde{\gamma})'(t)dt \\
        &= \int_a^b f(\gamma(\phi(t)))(\gamma'(\phi(t)))\phi'(t)dt \\
        &= \int_{a'}^{b'}f(\gamma(u))\gamma'(u)du \\
        &= \int_{\gamma}f(z)dz.
    \end{aligned}
    \]
\end{proof}

\begin{theorem}
\label{thm: cauchy-index-thm}
    Assume $\gamma$ is a closed path and $\Omega$ is the complement of $\gamma^*$, and 
    \[
    \text{Ind}_{\gamma}(z):= \dfrac{1}{2\pi i}\int_{\gamma}\dfrac{d\zeta}{\zeta - z},\quad z\in \Omega
    \]
    then $\text{Ind}_{\gamma}$ is a integer-valued function on $\Omega$ and it is a constant in any branch of $\Omega$. In particular, $\text{Ind}_{\gamma} = 0$ in the unbounded branch of $\Omega$.
\end{theorem}
\begin{proof}
    To show that
    \[
    \text{Ind}_{\gamma}(z) = \dfrac{1}{2\pi i}\int_a^b \dfrac{\gamma'(t)}{\gamma(t)-z}dt
    \]
    is integer-valued, notice that $e^{\omega}$ is $1$ iff $\omega$ is a integer multiple of $2\pi i$ and hence we consider
    \[
    \phi(t) = \exp\left(\int_a^t \dfrac{\gamma'(t)}{\gamma(t)-z}dt\right)
    \]
    where
    \[
    \phi'(t) = \phi(t)\left(\dfrac{\gamma'(t)}{\gamma(t)-z}\right) \implies \left(\dfrac{\phi}{\gamma-z}\right)'(t) = 0 .
    \]
    Therefore, we know $\phi(b)/(\gamma(b)-z) = \phi(a)/(\gamma(a) - z)$, and hence
    \[
    (2\pi i) \text{Ind}_{\gamma}(z) = \phi(b) = \phi(a) = 1
    \]
    since $\gamma(b) = \gamma(a)$, so $\text{Ind}_{\gamma}$ is always integer-valued. By the theorem \ref{thm: route-integral-holomorphic}, we know $\text{Ind}_{\gamma} \in H(\Omega)$ and hence it has to be constant on any connected branch by \ref{lem: connected-image}. For $z$ in the unbounded branch, since $|\gamma(t) - z|$ can be large enough such that $|\text{Ind}_{\gamma}(z)| < 1$, then it has to be $0$ and we are done. 
\end{proof}

\begin{proposition}
    If $\gamma$ is a positive-directed cycle with center $a$ and radius $r$, then
    \[
    \text{Ind}_{\gamma}(z) = \begin{cases}
        1,\quad |z-a| < r
        0,\quad |z-a| > r
    \end{cases}.
    \]
\end{proposition}
\begin{proof}
    We only need to prove the first case, that we know
    \[
    \text{Ind}_{\gamma}(a) = \int_0^{1} \dfrac{re^{2\pi i t}}{(re^{2\pi i t} +a) - a} dt =1 
    \]
    and we are done.
\end{proof}

\begin{theorem}
    Assume $F \in H(\Omega)$ and $F'$ continuous in $\Omega$, then for any closed path $\gamma$ in $\Omega$, then
    \[\int_{\gamma} F'(z)dz = 0.\]
\end{theorem}
\begin{proof}
    We know
    \[
    \int_{\gamma} F'(z)dz = \int_a^b F'(\gamma(z))\gamma'(z)dz = F(\gamma(b)) - F(\gamma(a)).
    \]
\end{proof}

\begin{corollary}
    
\end{corollary}
